\documentclass[11pt]{article}

%------------------------- fonts & encoding -----------------------------------
\usepackage[T1]{fontenc}
\usepackage[utf8]{inputenc}
\usepackage{amsmath,amssymb,amsthm}
\usepackage{mathtools}
\usepackage[sc]{mathpazo}
\linespread{1.06}
\usepackage{microtype}

%------------------------- layout ---------------------------------------------
\usepackage[a4paper,margin=1.05in,marginparwidth=1.0in]{geometry}
\usepackage{booktabs}
\usepackage{enumitem}
\setlist{itemsep=2pt,topsep=3pt}
\usepackage{xcolor}
\usepackage{graphicx}
\usepackage{tikz}
\usepackage{pgfplots}
\pgfplotsset{compat=1.17}
\usepackage[most,breakable]{tcolorbox}

%------------------------- palette --------------------------------------------
\definecolor{crimson}{HTML}{A51C30}
\definecolor{ink}{HTML}{1A1A1A}
\definecolor{steel}{HTML}{2C5F8A}
\definecolor{amber}{HTML}{9A6A00}
\definecolor{teal}{HTML}{1F6F6B}
\definecolor{paperlt}{HTML}{F7F4EF}

%------------------------- section styling ------------------------------------
\usepackage{titlesec}
\titleformat{\section}{\Large\bfseries\color{steel}}{\thesection}{0.7em}{}
\titleformat{\subsection}{\large\bfseries\color{ink}}{\thesubsection}{0.6em}{}
\titlespacing*{\section}{0pt}{1.6em}{0.6em}

%------------------------- hyperref -------------------------------------------
\usepackage[colorlinks=true,linkcolor=steel,citecolor=steel,urlcolor=crimson]{hyperref}

%------------------------- theorem environments -------------------------------
\theoremstyle{definition}
\newtheorem{definition}{Definition}[section]
\newtheorem{example}{Example}[section]
\newtheorem{exercise}{Exercise}[section]
\theoremstyle{plain}
\newtheorem{proposition}{Proposition}[section]
\theoremstyle{remark}
\newtheorem*{remark}{Remark}

%------------------------- callout boxes --------------------------------------
\newtcolorbox{intuition}{
  enhanced, breakable, colback=crimson!5, colframe=crimson, boxrule=0.5pt,
  left=8pt,right=8pt,top=6pt,bottom=6pt, arc=2pt,
  colbacktitle=crimson, coltitle=paperlt, fonttitle=\bfseries, title={Intuition}}

\newtcolorbox{pitfall}{
  enhanced, breakable, colback=amber!7, colframe=amber, boxrule=0.5pt,
  left=8pt,right=8pt,top=6pt,bottom=6pt, arc=2pt,
  colbacktitle=amber, coltitle=paperlt, fonttitle=\bfseries, title={Pitfall}}

\newtcolorbox{keyresult}{
  enhanced, breakable, colback=steel!6, colframe=steel, boxrule=0.5pt,
  left=8pt,right=8pt,top=6pt,bottom=6pt, arc=2pt,
  colbacktitle=steel, coltitle=paperlt, fonttitle=\bfseries, title={Takeaway}}

\newtcolorbox{observed}{
  enhanced, breakable, colback=teal!6, colframe=teal, boxrule=0.5pt,
  left=8pt,right=8pt,top=6pt,bottom=6pt, arc=2pt,
  colbacktitle=teal, coltitle=paperlt, fonttitle=\bfseries, title={Measured}}

%------------------------- math operators -------------------------------------
\DeclareMathOperator{\Var}{Var}
\DeclareMathOperator{\Cov}{Cov}
\DeclareMathOperator{\Corr}{corr}
\DeclareMathOperator{\diag}{diag}
\DeclareMathOperator{\std}{std}
\DeclareMathOperator*{\avg}{\mathbb{E}}
\newcommand{\T}{^{\!\top}}
\newcommand{\R}{\mathbb{R}}
\newcommand{\half}{\tfrac12}
\newcommand{\wvec}{\mathbf{w}}
\newcommand{\Xmat}{\mathbf{X}}
\newcommand{\Fmat}{\mathbf{F}}
\newcommand{\Vmat}{\mathbf{V}}
\newcommand{\Umat}{\mathbf{U}}
\newcommand{\Rmat}{\mathbf{R}}
\newcommand{\Lam}{\boldsymbol{\Lambda}}
\newcommand{\Del}{\boldsymbol{\Delta}}
\newcommand{\fvec}{\mathbf{f}}

%------------------------- title ----------------------------------------------
\title{\vspace{-2.2em}
  {\color{steel}\rule{\textwidth}{1.3pt}}\\[0.7em]
  \Huge\bfseries The Cross-Sectional Regression\\[0.35em]
  \large\mdseries Where factor returns come from --- and the country factor at its core\\[0.4em]
  {\color{steel}\rule{\textwidth}{0.7pt}}}
\author{}
\date{}

%------------------------- watermark footer -----------------------------------
\usepackage{fancyhdr}
\pagestyle{fancy}
\fancyhf{}
\renewcommand{\headrulewidth}{0pt}
\renewcommand{\footrulewidth}{0pt}
\fancyfoot[L]{\footnotesize\color{black!45}MSCI Barra USE4 Imitation}
\fancyfoot[R]{\footnotesize\color{black!45}\thepage}

\begin{document}
\maketitle
\thispagestyle{empty}

\begin{abstract}
Every factor so far has been an \emph{exposure} -- a number describing each stock. This note is where those exposures finally meet returns. Once a month the model runs a single cross-sectional regression: it regresses the realised returns of all stocks on their exposures, and the slopes that come out are the \emph{factor returns} -- what the country, each industry, and each style actually earned that month -- while the leftover, the part of each stock's return no factor explains, is its \emph{specific return}. This one regression is the engine of the model: it turns a wall of exposures into the two ingredients every risk forecast is built from. Three design choices make it work, and we take each in turn: the country factor enters as an all-ones intercept and turns out to be the market itself; the regression is weighted by the square root of market cap -- the inverse-variance optimum, which also happens to sit between equal- and cap-weighting; and the industry factors, which would otherwise be collinear with the country factor, are pinned down by a cap-weighted-sum-zero constraint. The note covers the country anchor and the regression together because they are one mechanism -- the country factor is the regression's intercept, and its only job is to anchor the very system this stage solves.
\end{abstract}

\begin{tcolorbox}[colback=paperlt,colframe=ink!30,boxrule=0.4pt,arc=2pt,left=10pt,right=10pt,top=8pt,bottom=8pt]
\textbf{How to read these notes.} We move from \emph{why} a cross-sectional regression is the source of factor returns, to the three things that make it well posed -- the country intercept (Section~\ref{sec:country}), the square-root-cap weighting (Section~\ref{sec:wls}), and the identification constraint (Section~\ref{sec:constraint}) -- and then to what it produces and what it explains. Read Sections~\ref{sec:wls} and~\ref{sec:constraint} slowly; they carry the weight of the set. The numbers labelled as measured come from a single run of our personal build.

\smallskip
Four coloured boxes recur: \textcolor{crimson}{\textbf{Intuition}} gives the mental picture, \textcolor{amber}{\textbf{Pitfall}} flags what breaks without the fix, \textcolor{steel}{\textbf{Takeaway}} states a load-bearing result, and \textcolor{teal}{\textbf{Measured}} reports real numbers from the build run.

\smallskip
\textbf{Reference material.}
\begin{itemize}[leftmargin=1.4em,topsep=2pt]
  \item External: Menchero, Orr \& Wang (2011), USE4 Methodology Notes (Sec. 2.1) and Empirical Notes (Sec. 4).
\end{itemize}
\end{tcolorbox}

\tableofcontents
\bigskip

\section{What the regression is, and why it is the engine}\label{sec:what}

\paragraph{Learning objectives.}
\begin{itemize}[leftmargin=1.4em,topsep=2pt]
  \item State the cross-sectional regression as returns-on-exposures, once per month.
  \item Distinguish the factor returns it estimates from the specific returns it leaves behind.
\end{itemize}

Up to now every factor has been an \emph{exposure}: a column describing each stock -- how large, how cheap, how volatile, which industry. Exposures alone forecast nothing; they have to be connected to returns. The cross-sectional regression is that connection. At each month-end the model takes the realised return of every stock over the coming month and regresses it, across the whole universe at once, on that month's exposures. The fitted slopes are the \emph{factor returns}: the single number each factor earned that month -- what it paid to be a bank, to be cheap, to be large. The residual that the regression cannot fit -- the part of a stock's return unique to it -- is its \emph{specific return}.

Written compactly, for the cross-section of stocks in a given month,
\[
y \;=\; \Xmat\,\fvec \;+\; u,
\]
where $y$ stacks the stocks' realised returns, $\Xmat$ is the exposure matrix, $\fvec$ is the vector of factor returns the regression solves for, and $u$ is the vector of specific returns. This is the inverse of the picture the overview painted: there, factor returns combine with exposures to make stock returns; here, we observe stock returns and exposures and \emph{solve} for the factor returns. One regression a month, repeated across the history, produces the entire factor-return record -- and that record, together with the specific returns, is what every covariance and risk forecast downstream is built from.

\begin{intuition}
Think of one month as a giant cross-sectional scatter: each stock is a point, its return on one axis and its exposures as coordinates. The regression asks, ``given how every stock was positioned, what single move best explains the spread of returns?'' The answer -- one number per factor -- is the factor return. Whatever distance remains between a stock and the fitted surface is its own private story: its specific return.
\end{intuition}

\begin{keyresult}
The cross-sectional regression $y=\Xmat\fvec+u$, run monthly, is where factor returns come from. It is the single hinge between the exposure layers (which describe stocks) and the risk model (which forecasts them), and it yields two deliverables at once: the factor returns $\fvec$ and the specific returns $u$.
\end{keyresult}

\section*{Notation}
\addcontentsline{toc}{section}{Notation}

\begin{center}
\begin{tabular}{@{}ll@{}}
\toprule
Symbol & Meaning \\
\midrule
$y_i$ & realised excess (log) return of stock $i$ over the month \\
$\Xmat=[\mathbf{1}\mid D\mid S]$ & exposure matrix: country intercept, $55$ industry dummies $D$, $12$ style scores $S$ \\
$\fvec$ & factor-return vector ($68 = 1+55+12$): $f_c$ (country), $f^{\text{ind}}_j$, $f^{\text{sty}}_k$ \\
$u_i$ & specific (idiosyncratic) return of stock $i$: $u_i = y_i - \Xmat_i\fvec$ \\
$v_i$ & regression weight $\propto \sqrt{\text{mcap}_i}$, normalised to sum to $1$ per date \\
$w_j$ & industry $j$'s cap weight; constraint $\sum_j w_j f^{\text{ind}}_j = 0$ \\
$R^2$ & weighted cross-sectional coefficient of determination for the month \\
ESTU & estimation universe (the regression sample; MSCI USA IMI proxy) \\
\bottomrule
\end{tabular}
\end{center}

\section{The country factor: the anchor at the intercept}\label{sec:country}

\paragraph{Learning objectives.}
\begin{itemize}[leftmargin=1.4em,topsep=2pt]
  \item State the country exposure and the role it plays in the regression.
  \item Explain why the country factor turns out to be the market.
\end{itemize}

The first column of $\Xmat$ is the country factor, and it is the simplest exposure in the whole model: every stock loads on it with value exactly $1$. It is the regression's intercept -- the move common to everything. In a single-country model like this one it is, in effect, the market factor; the name ``country'' is the multi-region framework's vocabulary, but here it is one number for the United States, and every stock is fully exposed to it.

Because the country factor is the part of return shared by all stocks, its estimated return $f_c$ comes out as essentially the market's return. This is not imposed -- it is a consequence of the constraint in Section~\ref{sec:constraint}, which routes the market-wide move into the country factor and leaves the industries and styles as deviations from it. Empirically the match is near-exact (Section~\ref{sec:outputs}): the country factor return tracks the cap-weighted market month for month. The country factor's only standalone deliverable, then, is not a return at all but an \emph{anchor} for the regression: the all-ones exposure column and the regression weights that the next section explains.

\begin{intuition}
Every stock rises a bit when the market rises; that shared lift has to be booked somewhere, or it would be double-counted across the industries. The country intercept is the account it is booked to. Make every stock load $1$ on it, force the industries to net to zero around it, and the intercept automatically becomes ``what the whole market did'' -- leaving each other factor to describe only how its members differed from that.
\end{intuition}

\section{Weighted least squares: why the square root of cap}\label{sec:wls}

\paragraph{Learning objectives.}
\begin{itemize}[leftmargin=1.4em,topsep=2pt]
  \item Explain why the regression is weighted at all, and why neither equal- nor cap-weighting is right.
  \item State what $\sqrt{\text{mcap}}$ weighting achieves.
\end{itemize}

The regression is not ordinary least squares; it is \emph{weighted}, and the choice of weights is a real modelling decision. Ordinary least squares would treat every stock as an equally reliable observation, but they are not: a tiny micro-cap's monthly return is mostly idiosyncratic noise, while a megacap's is steadier. Standard regression theory says to weight each observation by the inverse of its noise variance. The model takes a stock's specific-return variance to fall with size -- specifically, inversely with the \emph{square root} of market cap -- so inverse-variance weighting lands on a weight proportional to $\sqrt{\text{mcap}}$. That choice avoids both extremes at once: full cap-weighting would let a handful of giants dominate the fit so completely that the factor returns describe little more than the largest names, while equal weighting would hand the estimate to a sea of noisy small stocks.

So the square root of market cap is not an arbitrary compromise -- it is the inverse-variance (generalised-least-squares) \emph{optimum} under that variance model. Weighting each stock by $v_i \propto \sqrt{\text{mcap}_i}$, normalised to sum to one each month, damps the influence of the largest names without surrendering the regression to the smallest, and it also sits cleanly between the two extremes on a concentration curve. Figure~\ref{fig:country_dist} shows the three schemes: cap-weighting piles almost half the regression's attention onto the top one percent of names, equal weighting spreads it flat, and the square-root rule sits cleanly between -- the top one percent carry about a tenth.

\begin{figure}[htbp]
\centering
\includegraphics[width=\linewidth]{figures/country_dist.png}
\caption{How much of the regression's weight each scheme places on the largest stocks, as a concentration curve over one in-ESTU cross-section (stocks sorted small to large). Cap weighting (dashed) is extreme -- the top $1\%$ of names carry $\approx 48\%$ of the weight -- while equal weighting (dotted) is the diagonal. The square-root-of-cap weighting the model uses (solid) sits between them: the top $1\%$ carry $\approx 11\%$. It down-weights noisy micro-caps without letting a few giants dictate the factor returns.}
\label{fig:country_dist}
\end{figure}

\begin{pitfall}
The weighting is not cosmetic. Run the regression equal-weighted and a thousand tiny, noisy stocks outvote the market, so the factor returns wander on micro-cap noise. Run it fully cap-weighted and the few largest names \emph{are} the regression, so a factor's return reflects only how its mega-cap members did. Either extreme produces factor returns that are technically valid and practically useless; the square-root weight is what keeps the estimate both stable and representative.
\end{pitfall}

\section{The identification constraint}\label{sec:constraint}

\paragraph{Learning objectives.}
\begin{itemize}[leftmargin=1.4em,topsep=2pt]
  \item Explain why the industry columns and the country intercept are collinear.
  \item State the constraint that resolves it and how the build imposes it.
\end{itemize}

There is a redundancy built into $\Xmat$ that, left alone, makes the regression unsolvable. Every stock belongs to exactly one industry, so the $55$ industry dummy columns sum, row by row, to a column of all ones -- which is precisely the country column. The country factor is therefore an exact linear combination of the industries: the design matrix is rank-deficient, and infinitely many $(f_c, f^{\text{ind}})$ combinations fit the data equally well (add a constant to every industry return, subtract it from the country return, and nothing changes). The regression cannot choose among them without help.

The help is the identification constraint: the cap-weighted industry returns are required to sum to zero,
\[
\sum_{j} w_j\, f^{\text{ind}}_j \;=\; 0 ,
\]
with $w_j$ the industry cap weights built alongside the industry layer. This single equation removes exactly the one degree of freedom that was loose, and it does so meaningfully: it declares that the market-wide move lives in the country factor and the industries describe only deviations around it. The build imposes it exactly, not approximately -- it drops one industry column from the regression (the largest by cap weight, for the best conditioning), solves the smaller, now full-rank system, and then recovers the dropped industry's return from the constraint, $f_e = -\big(\sum_{j\neq e} w_j f_j\big)/w_e$. The choice of which industry to drop changes the arithmetic's conditioning but not the answer; this is exact restricted least squares.

\begin{intuition}
Without the constraint the model literally cannot tell ``the market rose $1\%$'' from ``every industry rose $1\%$'' -- they are the same fitted returns wearing different labels, and the regression has no reason to prefer one story. Forcing the cap-weighted industry returns to sum to zero settles it by convention: the common part is the country factor's, the industries only say who beat the market and who lagged. It is a choice of bookkeeping, but a necessary one, and it is what makes the country factor come out as the market.
\end{intuition}

\begin{keyresult}
The country and industry columns are collinear because every stock is in one industry. The cap-weighted-sum-zero constraint breaks the tie exactly -- the build eliminates one industry and recovers it -- giving the country factor the market and each industry a pure relative return.
\end{keyresult}

\section{Two outputs: factor returns and specific returns}\label{sec:outputs}

\paragraph{Learning objectives.}
\begin{itemize}[leftmargin=1.4em,topsep=2pt]
  \item Name the two deliverables and what each feeds downstream.
  \item Explain the in-sample versus out-of-sample treatment of specific returns.
\end{itemize}

Each solved month yields two things. The first is the vector of $68$ \emph{factor returns} -- one for the country, one for each of the $55$ industries, one for each of the $12$ styles -- the monthly record from which the factor covariance matrix is later estimated. The second is the \emph{specific return} of every stock, $u_i = y_i - \Xmat_i\fvec$: the part of its return no factor explains, the raw material of the specific-risk model. Two identities hold by construction and are checked: each stock's return splits exactly, $y_i = \hat y_i + u_i$, and the weighted specific returns sum to zero across the estimation universe each month, so the specific returns carry no leftover market exposure.

The specific returns are extended beyond the regression sample. The estimation universe is the liquid core on which the factors are \emph{fit}, but the model must forecast risk for the wider coverage universe too. So the fitted factor returns are applied \emph{out of sample} to the non-ESTU coverage stocks -- their specific return is computed from the same $\fvec$, with no re-estimation -- and flagged as out-of-sample. Those residuals are wider than the in-sample ones -- chiefly because the non-ESTU coverage names are smaller and less liquid, and so genuinely more idiosyncratic, with their out-of-sample treatment (no fit to them, no orthogonality guarantee) compounding the gap. That they run wider is exactly as it should be.

\begin{observed}
On the build run the regression solves $303$ monthly transitions (the tail dates skipped once the excess returns run out at the Ken French risk-free series' last date, which bounds the volatility factors, and the earliest dates skipped where a whole exposure column was not yet populated), producing $20{,}604$ factor-return rows ($303$ country $+$ $16{,}665$ industry $+$ $3{,}636$ style; a few dozen industry-months carry a null factor return, where an industry had no estimable members that month) and $\approx 1.10$M specific-return stock-dates. The in-sample (ESTU) specific returns have a standard deviation of $\approx 0.094$; the out-of-sample (non-ESTU) ones run wider at $\approx 0.175$, as expected. The two identities hold to machine precision: $y=\hat y+u$ exactly, and the weighted ESTU specific return is zero to $\approx 10^{-16}$.
\end{observed}

\section{What the regression recovers}\label{sec:diag}

\paragraph{Learning objectives.}
\begin{itemize}[leftmargin=1.4em,topsep=2pt]
  \item Read the cross-sectional $R^2$ and what it says about factor versus specific risk.
  \item Cite the evidence that the country factor is the market and the system is well conditioned.
\end{itemize}

How much of the cross-section do the factors actually explain? Month by month, a weighted cross-sectional $R^2$ of about a quarter (Figure~\ref{fig:csr_r2}): the $68$ factors together account for roughly $25\%$ of the spread in stock returns, and the remaining three-quarters is specific. That split is the whole reason the model has two halves -- a factor part and a specific part -- and it is why specific risk is not an afterthought but the larger share of a typical stock's variance. The $R^2$ is not constant: it spikes in crises, when the market moves as one and the country factor explains almost everything, and subsides in calm, stock-picker's markets where idiosyncratic moves dominate.

\begin{figure}[htbp]
\centering
\includegraphics[width=\linewidth]{figures/csr_r2.png}
\caption{The weighted cross-sectional $R^2$ of the monthly regression over the sample. It averages $\approx 25\%$ -- the factors explain about a quarter of the cross-section of returns, the rest being specific -- and spikes toward $50$--$75\%$ in crises (2008--09, 2020, 2022), when the market-wide move dominates and the country factor carries the day. The model's two halves, factor and specific risk, exist precisely because this number is well below one.}
\label{fig:csr_r2}
\end{figure}

\begin{observed}
On the build run the monthly cross-sectional $R^2$ averages $\approx 0.25$ (median $\approx 0.24$), ranging from $\approx 0.09$ in quiet months to $\approx 0.76$ in crisis months. The country factor return correlates $\approx 0.998$ with the cap-weighted market and has an annualised volatility of $\approx 16\%$ -- it is, to a very close approximation, the market. The identification constraint holds to $\approx 10^{-18}$, and the square-root-scaled design is well conditioned throughout, with a worst-case condition number of $\approx 50$ -- orders of magnitude \emph{below} the ill-conditioning danger zone, the payoff of imposing the constraint by elimination rather than fighting the collinearity head-on.
\end{observed}

\section{Summary, and one caveat}\label{sec:summary}

\paragraph{Learning objectives.}
\begin{itemize}[leftmargin=1.4em,topsep=2pt]
  \item Recite the monthly regression from exposures to the two deliverables.
  \item State the central limitation and what it means downstream.
\end{itemize}

The stage, end to end: each month assemble the exposure matrix $\Xmat=[\mathbf{1}\mid D\mid S]$ from the country anchor, the industry dummies, and the style scores; weight every stock by the square root of its market cap; regress realised returns on the exposures under the cap-weighted-sum-zero industry constraint, imposed exactly by eliminating one industry and recovering it; and ship the $68$ factor returns and the per-stock specific returns, the latter extended out of sample to the full coverage universe. The country factor is the all-ones intercept of this regression, and it emerges as the market.

\begin{keyresult}
The cross-sectional regression is the model's engine: it converts a month of exposures and returns into factor returns and specific returns. Its three load-bearing choices are the country intercept (the market anchor), the square-root-cap weighting (stable and representative), and the identification constraint (which makes the system solvable and the country factor the market).
\end{keyresult}

\begin{pitfall}
One honest limitation is structural to the method: the factor returns are \emph{estimates} from a single monthly cross-section, not observed quantities, and they are only as good as that month's sample and exposures. A month with thin coverage of some exposure, or with too few stocks, cannot be solved cleanly and is skipped -- so the factor-return history has gaps at the sparse early dates and the data's tail. And because only about a quarter of the cross-section is explained, the factor returns are recovered against a large cloud of specific noise; they are stable in aggregate but should never be read as precise monthly truths. Downstream, both the factor covariance and the specific-risk model must treat their inputs as noisy estimates, which is exactly why those stages lean so heavily on smoothing and shrinkage rather than trusting any single month.
\end{pitfall}

\begin{remark}
None of this changes the stage's place in the model: the cross-sectional regression is the hinge on which everything turns -- exposures in, factor and specific returns out. The caveats are about \emph{estimation noise in a single month}, not about whether the decomposition into common and specific return is sound; that decomposition is the foundation the rest of the model stands on.
\end{remark}

\begin{center}\small\itshape
References: Menchero, Orr \& Wang (2011), \textit{The Barra US Equity Model (USE4)} --- Methodology Notes (Sec. 2.1) and Empirical Notes (Sec. 4). ESTU treated as an MSCI USA IMI proxy; measured quantities are from a personal build run.
\end{center}

\end{document}
